% GENERATED FILE. Edit canonical lessons, then run npm run generate:books.
% canonical-source-hash: fnv1a64:d3f762cec1caca90
% canonical-lessons: ES-C651-vagon, ES-C651-maquinista, ES-C651-taquilla, ES-C651-litera, ES-C651-apearse

\chapter{En el tren}
\label{ch:es-en-el-tren}
\hlchaptermodality{651}

\begin{chapteropening}
\textbf{By the end of this chapter:} \emph{I can travel by train: find my carriage, sleep in a berth, buy a ticket at the ticket office and get off at my stop.}

Complete the last lesson of chapter 651: I can travel by train: find my carriage, sleep in a berth, buy a ticket at the ticket office and get off at my stop.
\end{chapteropening}

\section[el vagón]{el vagón — the carriage, the railway car}

\label{lesson:ES-C651-vagon}



\begin{quote}
Before the new one: say the Spanish for the change, then the Spanish for to validate.
\end{quote}

\subsection*{You'll want to know: el vagón}
\textbf{el vagón} — \textquotedblleft{}the carriage, the railway car\textquotedblright{}.

\begin{quote}
Nuestro asiento está en el vagón tres.
\end{quote}

\begin{quote}
\emph{Our seat is in carriage three.}
\end{quote}

\emph{El vagón restaurante} — the dining car.

\begin{cousinweb}
\emph{vagón} was borrowed from English \emph{wagon} (EtymDB), which itself came from Dutch \emph{wagen}.
\end{cousinweb}

\begin{grammarlens}[title={What you've built}]
One more word for telling, explaining and describing at B1.
\end{grammarlens}

\subsection*{Your turn}
\begin{itemize}
\raggedright
  \item \emph{Say it:} \emph{el vagón}
  \item \emph{Say it:} \emph{el vagón}, once more
  \item \emph{Say it:} a sentence of your own with \emph{el vagón}
  \item \emph{Recall:} say the Spanish for the change, then the Spanish for to validate, then say \emph{el vagón} again
\end{itemize}

\begin{tcolorbox}[breakable,colback=teal!4,colframe=teal!35!black,title={Before you move on}]
What does el vagón mean? (\textquotedblleft{}The carriage, the railway car\textquotedblright{}.) Say it once more.
\end{tcolorbox}
\section[el maquinista]{el maquinista — the train driver, the engine driver}

\label{lesson:ES-C651-maquinista}



\begin{quote}
Before the new one: say the Spanish for to validate, then the Spanish for the carriage.
\end{quote}

\subsection*{You'll want to know: el maquinista}
\textbf{el maquinista} — \textquotedblleft{}the train driver, the engine driver\textquotedblright{}.

\begin{quote}
El maquinista para el tren en la estación.
\end{quote}

\begin{quote}
\emph{The train driver stops the train at the station.}
\end{quote}

Built on \emph{la máquina}: the first trains were pulled by steam engines.

\begin{cousinweb}
\emph{maquinista} is built on \emph{máquina} \textquotedblleft{}machine\textquotedblright{} with \emph{-ista} (EtymDB); \emph{máquina} is a learned borrowing from Latin \emph{māchina}, from Greek \emph{mēkhan\'{\={e}}}.
\end{cousinweb}

\begin{grammarlens}[title={What you've built}]
One more word for telling, explaining and describing at B1.
\end{grammarlens}

\subsection*{Your turn}
\begin{itemize}
\raggedright
  \item \emph{Say it:} \emph{el maquinista}
  \item \emph{Say it:} \emph{el maquinista}, once more
  \item \emph{Say it:} a sentence of your own with \emph{el maquinista}
  \item \emph{Recall:} say the Spanish for to validate, then the Spanish for the carriage, then say \emph{el maquinista} again
\end{itemize}

\begin{tcolorbox}[breakable,colback=teal!4,colframe=teal!35!black,title={Before you move on}]
What does el maquinista mean? (\textquotedblleft{}The train driver, the engine driver\textquotedblright{}.) Say it once more.
\end{tcolorbox}
\section[la taquilla]{la taquilla — the ticket office, the box office; the locker}

\label{lesson:ES-C651-taquilla}



\begin{quote}
Before the new one: say the Spanish for the carriage, then the Spanish for the train driver.
\end{quote}

\subsection*{You'll want to know: la taquilla}
\textbf{la taquilla} — \textquotedblleft{}the ticket office, the box office; the locker\textquotedblright{}.

\begin{quote}
Compramos los billetes en la taquilla.
\end{quote}

\begin{quote}
\emph{We bought the tickets at the ticket office.}
\end{quote}

In Latin America also \emph{la boletería}.

\begin{cousinweb}
\emph{taquilla} is a little \emph{taca} \textquotedblleft{}a small cupboard\textquotedblright{}. Where \emph{taca} comes from is not settled.
\end{cousinweb}

\begin{grammarlens}[title={What you've built}]
One more word for telling, explaining and describing at B1.
\end{grammarlens}

\subsection*{Your turn}
\begin{itemize}
\raggedright
  \item \emph{Say it:} \emph{la taquilla}
  \item \emph{Say it:} \emph{la taquilla}, once more
  \item \emph{Say it:} a sentence of your own with \emph{la taquilla}
  \item \emph{Recall:} say the Spanish for the carriage, then the Spanish for the train driver, then say \emph{la taquilla} again
\end{itemize}

\begin{tcolorbox}[breakable,colback=teal!4,colframe=teal!35!black,title={Before you move on}]
What does la taquilla mean? (\textquotedblleft{}The ticket office, the box office; the locker\textquotedblright{}.) Say it once more.
\end{tcolorbox}
\section[la litera]{la litera — the berth, the bunk}

\label{lesson:ES-C651-litera}



\begin{quote}
Before the new one: say the Spanish for the train driver, then the Spanish for the ticket office.
\end{quote}

\subsection*{You'll want to know: la litera}
\textbf{la litera} — \textquotedblleft{}the berth, the bunk\textquotedblright{}.

\begin{quote}
Dormimos en litera en el tren de noche.
\end{quote}

\begin{quote}
\emph{We slept in a berth on the night train.}
\end{quote}

At home, \emph{una litera} is a bunk bed.

\begin{cousinweb}
\emph{litera} came from Catalan \emph{llitera} (EtymDB), built on \emph{llit} \textquotedblleft{}bed\textquotedblright{}, which comes from Latin \emph{lectus} \textquotedblleft{}a couch, a bed\textquotedblright{} (Lewis \& Short). French \emph{litière} is a cousin.
\end{cousinweb}

\begin{grammarlens}[title={What you've built}]
One more word for telling, explaining and describing at B1.
\end{grammarlens}

\subsection*{Your turn}
\begin{itemize}
\raggedright
  \item \emph{Say it:} \emph{la litera}
  \item \emph{Say it:} \emph{la litera}, once more
  \item \emph{Say it:} a sentence of your own with \emph{la litera}
  \item \emph{Recall:} say the Spanish for the train driver, then the Spanish for the ticket office, then say \emph{la litera} again
\end{itemize}

\begin{tcolorbox}[breakable,colback=teal!4,colframe=teal!35!black,title={Before you move on}]
What does la litera mean? (\textquotedblleft{}The berth, the bunk\textquotedblright{}.) Say it once more.
\end{tcolorbox}
\section[apearse]{apearse — to get off (a train or bus), to step down}

\label{lesson:ES-C651-apearse}



\begin{quote}
Before the new one: say the Spanish for the ticket office, then the Spanish for the berth.
\end{quote}

\subsection*{You'll want to know: apearse}
\textbf{apearse} — \textquotedblleft{}to get off (a train or bus), to step down\textquotedblright{}.

\begin{quote}
Nos apeamos en la próxima estación.
\end{quote}

\begin{quote}
\emph{We are getting off at the next station.}
\end{quote}

A formal word heard in announcements; in speech people say \emph{bajarse}.

\begin{cousinweb}
\emph{apearse} is usually traced to a Vulgar Latin verb \emph{appedāre} \textquotedblleft{}to set foot down\textquotedblright{}, built on Latin \emph{pēs, pedis} \textquotedblleft{}a foot\textquotedblright{} (Lewis \& Short).
\end{cousinweb}

\begin{grammarlens}[title={What you've built}]
One more word for telling, explaining and describing at B1.
\end{grammarlens}

\subsection*{Your turn}
\begin{itemize}
\raggedright
  \item \emph{Say it:} \emph{apearse}
  \item \emph{Say it:} \emph{apearse}, once more
  \item \emph{Say it:} a sentence of your own with \emph{apearse}
  \item \emph{Recall:} say the Spanish for the ticket office, then the Spanish for the berth, then say \emph{apearse} again
\end{itemize}

\begin{tcolorbox}[breakable,colback=teal!4,colframe=teal!35!black,title={Before you move on}]
What does apearse mean? (\textquotedblleft{}To get off (a train or bus), to step down\textquotedblright{}.) Say it once more.
\end{tcolorbox}
